% GitHub Repo and Documentation: https://github.com/celiobjunior/resume-template
% Copyright © 2025 Celio B Junior. All rights reserved.
%
% Licensed under the Apache License, Version 2.0 (the "License");
% you may not use this file except in compliance with the License.
% You may obtain a copy of the License at
%
%     http://www.apache.org/licenses/LICENSE-2.0
%
% This template follows best practices from README.md

% Start of the LaTeX document: defines document type and format
% a4paper = A4 size, 10pt = base font size
\documentclass[a4paper,10pt]{article}

% --- PACKAGES: BASE / LANGUAGE ---
\usepackage[utf8]{inputenc}
\usepackage[T1]{fontenc}
\usepackage[english]{babel}

% --- PACKAGES: LAYOUT / TYPOGRAPHY ---
\usepackage{geometry}
\usepackage{parskip}
\usepackage{microtype}
\usepackage{enumitem}
\usepackage{titlesec}
\usepackage{array}
\usepackage{tabularx}

% --- PACKAGES: LINKS / BOOKMARKS ---
\usepackage{hyperref}
\usepackage{bookmark}
\usepackage{xurl}

% --- SETTINGS: MARGINS / SPACING ---
\geometry{top=1.5cm, bottom=1.5cm, left=1.5cm, right=1.5cm} % Page margins
\setcounter{secnumdepth}{0}                                 % Removes section numbering
\setlist[itemize]{
    leftmargin=0.75em, % List indentation (closer to page margin)
    itemsep=0.2em,     % Space between items
    topsep=0.25em,     % Space before/after list
    parsep=0em,        % Space between paragraphs in same item
    partopsep=0em      % Extra spacing at list start
}

% Removes page numbers and headers for a clean resume layout
\pagestyle{empty}

% --- SETTINGS: METADATA / LINKS ---
\hypersetup{
    pdftitle={Karolayne Amabile Borges},   % PDF title
    pdfauthor={Karolayne Amabile Borges},  % PDF author
    colorlinks=true,                       % Uses colored links (no boxes)
    linkcolor=black,                       % Internal links color
    urlcolor=black,                        % URLs color
    citecolor=black,                       % Citations color
    bookmarksdepth=2                       % Sidebar: section and subsection
}

% --- SETTINGS: TITLES ---
\titleformat{\section}
{\Large\bfseries}
{}
{0em}
{}
[\titlerule\vspace{0.5ex}]
\titleformat{\subsection}
{\normalsize\bfseries}
{}
{0em}
{}
\titlespacing*{\subsection}{0pt}{0.35em}{0.15em}

% Macro for resume entries with PDF bookmark
\newcounter{cventry}
\newcommand{\cventry}[4]{%
    \refstepcounter{cventry}%
    \phantomsection%
    \pdfbookmark[2]{#1}{cventry-\thecventry}%
    \noindent\begin{tabularx}{\textwidth}{@{}>{\raggedright\arraybackslash}X >{\raggedleft\arraybackslash}X@{}}
    \textbf{#1} & #2 \\
    \textit{#3} & \textit{#4} \\
    \end{tabularx}
}

% --- START OF DOCUMENT ---
\begin{document}

% --- HEADER ---
\begin{center}
    {\LARGE \textbf{Karolayne Amabile Borges}}
    \\ [0.1cm]
    Anápolis, Brazil
    {\textbullet}
    \href{mailto:karolayneamabile@gmail.com}{karolayneamabile@gmail.com}
    {\textbullet}
    \href{https://linkedin.com/in/karolayneamabile}{linkedin.com/in/karolayneamabile}
    {\textbullet}
    \href{https://github.com/KarolayneAmabile}{github.com/KarolayneAmabile}
\end{center}

% --- SECTIONS ---

\section{Summary}

Junior DevOps Engineer with hands-on experience building and maintaining cloud-native infrastructure
on AWS and Kubernetes. Proven track record of automating deployments, implementing GitOps
workflows, and establishing end-to-end observability in production environments. AWS Certified
Cloud Practitioner. Currently pursuing a B.S. in Computer Science at IFG (Brazil).

% ------

\section{Experience}
    \cventry{Soliton}{Remote}{Junior DevOps Engineer}{April 2025 -- May 2026}
        \begin{itemize}
            \item Reduced OpenEDX platform deployment time by 5$\times$ by replacing Tutor's
            abstraction layer with a custom GitHub Actions pipeline (optimized cache + direct
            Dockerfile), ArgoCD, and ArgoCD Image Updater.
            \item Designed and managed multi-purpose Kubernetes clusters (security, monitoring,
            operations, and production), enforcing environment isolation and governance per
            namespace.
            \item Developed reusable Terraform/OpenTofu modules and Helm charts to standardize
            infrastructure provisioning across environments, reducing deployment inconsistencies
            and enforcing GitOps practices.
            \item Implemented a full-stack observability solution with Grafana, Prometheus, Loki,
            and Alloy, enabling real-time monitoring and alerting for cloud-native workloads.
            \item Automated secrets lifecycle management using HashiCorp Vault and External
            Secrets Operator, eliminating hardcoded credentials from all repositories.
            \item Built a RESTful API in Rust following hexagonal architecture, decoupling
            domain logic from application and infrastructure layers.
        \end{itemize}

% ------

\section{Skills}
    \begin{itemize}
        \item \textbf{Cloud \& IaC:} AWS, Terraform/OpenTofu, Terragrunt, Helmfile
        \item \textbf{Kubernetes \& GitOps:} Kubernetes (EKS, k0s), ArgoCD, Helm,
        Cilium, Cloudflared, cert-manager
        \item \textbf{CI/CD:} GitHub Actions, ArgoCD, ArgoCD Image Updater
        \item \textbf{Observability:} Grafana, Prometheus, Loki, Alloy
        \item \textbf{Storage \& Backup:} Longhorn, Rook-Ceph, Velero
        \item \textbf{Security \& Secrets:} HashiCorp Vault, Vaultwarden, External Secrets Operator
        \item \textbf{Programming:} Rust
        \item \textbf{Languages:} Portuguese (Native), English (Intermediate)
    \end{itemize}

% ------

\section{Certifications}
    \begin{itemize}
        \item \textbf{AWS Certified Cloud Practitioner} — Amazon Web Services
    \end{itemize}

% ------

\section{Education}
    \cventry{Instituto Federal de Goiás (IFG)}{Goiás, Brazil}
    {B.S. in Computer Science}{2024 -- 2027}
    \begin{itemize}
        \item Academic Score: 8.0/10.0 — 5th semester
        \item Relevant coursework: Computer Architecture, Operating Systems, Computer
        Networks, Software Engineering, Data Structures, Databases
    \end{itemize}

\end{document}
